% ============================================================
\section{Graph (Directed Graphs)}
A categoria de grafos direcionados é essencial para modelar redes, fluxos e transições de estado.
Um objeto nesta categoria é constituído por um conjunto de vértices e arestas com direcionalidade explícita.

\subsection{Objetos}
Em MATLAB, os grafos são naturalmente modelados pelas suas matrizes de adjacência.
O valor da matriz na posição $(i,j)$ indica a presença de uma aresta direcionada do nó $i$ para o nó $j$.

\begin{lstlisting}[caption={Representação de um Grafo Direcionado}]
% Matriz de adjacência (Vértices = 3, Arestas = 2)
n = 3;
M = [0 1 0; 0 0 0; 0 1 0];
\end{lstlisting}

\subsection{Morfismos}
Os morfismos são homomorfismos de grafos: funções entre vértices que mapeiam arestas válidas no domínio para arestas válidas no codomínio.
Se existe uma aresta $i \to j$, tem de existir uma aresta $f(i) \to f(j)$.

\begin{lstlisting}[caption={Verificação de Homomorfismo de Grafos}]
% f mapeia vértices do domínio para o codomínio
isHomomorphism = true;
[I, J] = find(M_domain);
for k = 1:length(I)
    if ~M_codomain(f(I(k)), f(J(k)))
        isHomomorphism = false;
        break;
    end
end
\end{lstlisting}